% !TeX program = xelatex
% ============================================================
% project_proposal/project_proposal.tex — «Описание проекта и заявка на
% выполнение» курсового проекта 1 курса магистратуры СПИ ФКН НИУ ВШЭ.
%
% Форма — шаблон семинара наставника (ред. 2.1) + требования памятки
% «КП и ВКР на СПИ» (ред. 3) и слайдов семинара:
%   * предварительная редакция (PRP, до 15 ноября вместе с заявлением):
%     обязательны разделы со звёздочкой — название, автор, руководитель,
%     предметная область и мотивация, обзор существующих работ, постановка
%     задачи (исследовательский вопрос, объект, предмет), цель, ожидаемые
%     результаты, источники;
%   * доработанная редакция (EPP, контрольная точка CP1 в январе): плюс
%     аннотация, термины и сокращения, план работ, ресурсы и риски.
% Редакция переключается в настройках проекта («Редакция заявки») — блоки
% \ifhseExtendedProposal ниже появляются и пропадают без правки текста.
%
% Пишите свой текст вместо жёлтых \hseFill{…}. Данные о себе, руководителе
% и названии задайте в настройках проекта — они подставятся сами.
% ============================================================
\documentclass[11pt,a4paper]{article}
% !TeX program = xelatex
% ============================================================
% project_proposal/preamble.tex — преамбула «Описания проекта и заявки
% на выполнение» курсового проекта магистратуры СПИ.
%
% Повторяет вёрстку шаблона семинара наставника (ред. 2.1): A4, поля
% 20/15/20/20 мм, синие заголовки разделов, таблица «Сводка». Документ
% подаётся в PDF вместе с заявлением о теме (SmartLMS, до 15 ноября).
% Язык разделов — язык работы (project.lang).
% ============================================================
\newcommand{\hseDefaultMono}{Consolas}
\input{../common/fonts}
((* if project.lang == 'en' *))
\usepackage[russian,english]{babel}
((* else *))
\usepackage[english,russian]{babel}
((* endif *))

\usepackage[a4paper,left=20mm,right=15mm,top=20mm,bottom=20mm]{geometry}
\usepackage{setspace}
\setstretch{1.15}
\usepackage{indentfirst}
\setlength{\parindent}{1.25cm}
\setlength{\parskip}{4pt}

\usepackage{microtype}
\usepackage{csquotes}
\usepackage{amsmath}
\usepackage{amssymb}
\usepackage{etoolbox}
\usepackage{graphicx}
\usepackage{xcolor}
\usepackage{array}
\usepackage{tabularx}
\usepackage{booktabs}
\usepackage{longtable}
\usepackage{enumitem}
\setlist{nosep, leftmargin=1.25cm}
\setlist[itemize,1]{label=--}
\usepackage{xurl}
\urlstyle{same}
\usepackage{hyperref}
\hypersetup{unicode=true,colorlinks=true,linkcolor=black,urlcolor=black,citecolor=black}
\usepackage{cite}
\usepackage{lastpage}
\usepackage{fancyhdr}

% ── Цвета и заголовки разделов (как в шаблоне семинара) ─────
\definecolor{hseprpnote}{RGB}{230,184,183}
\usepackage{titlesec}
\titleformat{\section}{\large\bfseries}{\thesection.}{0.4em}{}
\titlespacing*{\section}{0pt}{14pt}{6pt}
\titleformat{\subsection}{\normalsize\bfseries}{}{0em}{}
\titlespacing*{\subsection}{0pt}{10pt}{4pt}
% Обязательный раздел (поле со звёздочкой в шаблоне ОП).
\newcommand{\hseReq}{\textsuperscript{*}}

% ── Колонтитул: номер страницы и редакция формы ─────────────
\pagestyle{fancy}
\fancyhf{}
\renewcommand{\headrulewidth}{0pt}
\fancyfoot[C]{\thepage}
\fancyfoot[R]{\scriptsize ((* if project.lang == 'en' *))Form rev. 2.1 of 1.11.2023((* else *))Ред. 2.1 от 1.11.2023((* endif *))}

% Номера в списке источников с точкой.
\makeatletter
\renewcommand\@biblabel[1]{#1.}
\makeatother

\input{../common/meta}

\begin{document}

% ── Шапка формы ─────────────────────────────────────────────
\begin{center}
{\Large\bfseries Курсовой проект}\\[2pt]
{(1-й год обучения в магистратуре СПИ)}\\[10pt]
{\Large\bfseries Описание проекта и заявка на выполнение}\\[2pt]
{\ifhseExtendedProposal(доработанное)\else(предварительное)\fi}
\end{center}

% ── Сводка ──────────────────────────────────────────────────
\subsection*{Сводка\hseReq}
\noindent
\renewcommand{\arraystretch}{1.35}
\begin{tabularx}{\textwidth}{|>{\raggedright\arraybackslash\bfseries}p{0.3\textwidth}|>{\raggedright\arraybackslash}X|}
\hline
1. Название работы на русском языке & \hseTitleRu \\
\hline
2. Название работы на английском языке & \hseTitleEn \\
\hline
3. Язык выполнения работы & \hseWorkLang \\
\hline
4. ФИО студента (автора работы) & \hseAuthorName \\
\hline
5. ФИО научного руководителя, его учёное звание, степень и должность & \hseSupervisorFull
\ifhseHasCoSupervisor\newline Консультант: \hseCoSupervisorName\ifdefempty{\hseCoSupervisorDegree}{}{, \hseCoSupervisorDegree}\ifdefempty{\hseCoSupervisorTitle}{}{, \hseCoSupervisorTitle}\fi \\
\hline
6. Подпись научного руководителя и дата подписания &
\vspace{1.6cm}
% Подпись руководителя ставится в инструменте «Подписи» (PNG) или
% от руки; при открепленной ЭП (GnuPG) замените строку ниже на
% «Подписано с помощью открепленной электронной подписи».
\signdateblank \\
\hline
\end{tabularx}

% ============================================================
\ifhseExtendedProposal
\subsection*{Аннотация}
% Аннотация будущей работы (ГОСТ 7.9): 1–2 абзаца, 500–1500 знаков.
% Отвечает на вопрос «о чём работа?» — не пересказ результатов.
\hseFill{о чём работа: проблема, подход, ожидаемый результат — один-два абзаца}
\fi

% ============================================================
\setcounter{section}{6}
\section{Описание предметной области, актуальность и мотивация к~её разработке (исследованию)\hseReq}
% Предметная область: границы исследования, контекст. Актуальность: почему
% это важно сейчас и кому. Мотивация: какой пробел (research gap) закрываем
% и почему беретесь именно вы.
\hseFill{кратко опишите предметную область и границы работы}

\hseFill{почему тема актуальна именно сейчас: кому и какую пользу принесёт решение}

\hseFill{мотивация: какой пробел в существующих решениях закрывает работа}

\subsection*{Обзор существующих работ и решений\hseReq}
% Что уже известно по теме и почему этого недостаточно: аналоги, подходы,
% статьи. Каждая работа — со ссылкой на источник из списка ниже.
\hseFill{существующие подходы и решения со ссылками на источники и их ограничения}~\cite{denscombe2019}.

\subsection*{Постановка задачи\hseReq}
% Лаконично, но исчерпывающе: какую конкретно проблему решает работа.
\hseFill{какую конкретно проблему планируется решить в рамках работы}

\noindent\textbf{Исследовательский вопрос:} \hseFill{вопрос, на который отвечает работа}.

\noindent\textbf{Объект исследования} — \hseFill{широкая область, процесс или явление, на которое направлена работа}.

\noindent\textbf{Предмет исследования} — \hseFill{конкретная часть, свойство или аспект объекта, который изучается}.

% ============================================================
\section{Ключевые слова}
% Ключевые слова разделяются точкой с запятой (;).
\hseFill{ключевое слово; ключевое слово; ключевое слово}

% ============================================================
\section{Цель работы\hseReq}
% Одна цель. Формула: глагол действия (разработать, определить, оценить)
% + объект + на основе чего / с помощью чего. Критерии SMART.
\hseFill{разработать … для … на основе …}

% ============================================================
\section{Задачи работы}
% 3–7 шагов от анализа к реализации и оценке результатов. Задачи
% формируют основу календарного плана — не укрупняйте их.
\begin{enumerate}[label=\arabic*)]
  \item \hseFill{проанализировать существующие решения и подходы};
  \item \hseFill{сформулировать требования к решению};
  \item \hseFill{спроектировать архитектуру (метод, модель)};
  \item \hseFill{реализовать прототип};
  \item \hseFill{провести испытания и оценить результаты}.
\end{enumerate}

% ============================================================
\section{Ожидаемые результаты\hseReq}
% Результаты ИЗМЕРИМЫ: должно быть легко проверить, достигнуты ли они.
% Каждый результат отвечает на задачу; при необходимости укажите методику
% проверки (метрика, порог, способ измерения).
\begin{enumerate}[label=\arabic*)]
  \item \hseFill{результат} --- \hseFill{как проверить: метрика и порог};
  \item \hseFill{результат} --- \hseFill{как проверить}.
\end{enumerate}

% ============================================================
\section{Применяемые методы}
% Ожидаемые: анализ и синтез, сравнительный анализ, моделирование,
% эксперимент, тестирование и т.д.
\hseFill{ожидаемые методы}

% ============================================================
\section{Используемые инструменты}
% Ожидаемые: языки, фреймворки, СУБД, инфраструктура, системы управления
% проектом (GitLab ФКН, трекер задач).
\hseFill{ожидаемые инструменты}

% ============================================================
\subsection*{Дополнительные разделы (опционально)}
\ifhseExtendedProposal
\medskip\noindent\textbf{Термины и определения.}\par
\begin{description}[leftmargin=0pt, font=\bfseries]
  \item[\hseFill{термин}] --- \hseFill{определение}.
\end{description}

\medskip\noindent\textbf{Сокращения и обозначения.}\par
\begin{description}[leftmargin=0pt, font=\bfseries]
  \item[\hseFill{АББР}] --- \hseFill{расшифровка}.
\end{description}

\medskip\noindent\textbf{План работ.}\par
% Календарный план на основе задач из п. 10: этап, срок, результат.
\noindent
\begin{tabularx}{\textwidth}{|p{0.45\textwidth}|p{0.15\textwidth}|X|}
\hline
\textbf{Этап} & \textbf{Срок} & \textbf{Результат} \\
\hline
\hseFill{задача 1} & \hseFill{месяц} & \hseFill{артефакт} \\
\hline
\hseFill{задача 2} & \hseFill{месяц} & \hseFill{артефакт} \\
\hline
\end{tabularx}

\medskip\noindent\textbf{Ресурсы и риски.}\par
% Выполнимость: время, данные, доступы, вычислительные ресурсы; этические
% и правовые ограничения; что делать, если риск сработает.
\hseFill{ресурсы, ограничения и основные риски с мерами реагирования}
\else
\hseOptional{Здесь можно привести дополнительную информацию, необходимую для раскрытия идеи курсового проекта, например приложения.}
\fi

% ============================================================
% Список источников*: на каждый источник есть ссылка в тексте выше.
\renewcommand{\refname}{Список источников\hseReq}
\begin{thebibliography}{99}
  \bibitem{denscombe2019}
  Denscombe~M. Research Proposals: A Practical Guide. --- 2nd ed. --- London: Open University Press, 2019. --- 176~p.
\end{thebibliography}

\end{document}
